\documentclass[,
  addpoints,
  % answers,
  coursename={测试课程},
  examtype={A},
  at={期末},
  examdate={2026年1月21日},
]{fzuexam}

\usepackage{zhlipsum}

\begin{document}
\question[30] 单项选择题（15题，每题2分，请将答案写到题号后的空格中）

\droptotalpoints

\begin{parts}
  \part \fillin[A][1.5em] 这是第一道选择题，choices环境排版选项，每个选项一行
  \begin{choices}
    \choice 选项1
    \choice 选项2
    \choice 选项3
    \choice 选项4
    \choice 选项5
  \end{choices}
  \part \fillin[C][1.5em] 这是第二道选择题，oneparchoices环境排版选项，所有选项一行，需要插入空行进行换行

  \begin{oneparchoices}
    \choice 选项1
    \choice 选项2
    \choice 选项3
    \choice 选项4
    \choice 选项5
  \end{oneparchoices}
\end{parts}

\question 填空题（2题，每题2分）

\droptotalpoints

\begin{parts}
  \part[2] 自监督学习中，通过构造正负样本对来学习表征的方法称为 \fillin[对比] 学习。
  \part[2] 在轨迹数据预处理中，将原始GPS点序列转换为语义停留区域的过程称为 \fillin[停留点检测]。
\end{parts}


\question[30] 综合题

\droptotalpoints

需要明确小分的地方需要暂时使用\verb`\noaddpoints`关闭分数统计 

（未关闭统计示例）
% \noaddpoints
\begin{parts}
  \part[10] \zhlipsum[1]
  \part[10] \zhlipsum[1]
\end{parts}


\question[30] 综合题

\droptotalpoints

（关闭统计示例）
\noaddpoints
\begin{parts}
  \part[10] \zhlipsum[1]
  \part[10] \zhlipsum[1]
\end{parts}

\addpoints

\question[10] 简答题（2题，每题5分）

\droptotalpoints

\begin{parts}
  \part \zhlipsum[1]
  \begin{solutionorlines}[6cm]
    \zhlipsum[1]
  \end{solutionorlines}

  \part \zhlipsum[1]
  \begin{solutionorlines}[6cm]
    \zhlipsum[1]
  \end{solutionorlines}
\end{parts}

\end{document}